\begin{lemma}[Imported neural structural bounds and their budgets]
\label{lem:neural-inputs}
Fix depth $L\ge2$, let $C=C_L(f)<\infty$, and suppose $N\ge C$.
All function priors in this lemma use the global output margin restriction
before training. There are functions $F_L(d,J)$, $A_L(d,C)\ge1$,
$T_L(d,B)$, and $G_L(d,B)$ with polynomial majorants at fixed $L$
such that the following statements hold. The first is the norm count
\[
\log\#\{g\in\mathcal H_N:\kappa_N(g)\le J\}\le F_L(d,J)
 \quad(J\ge1),
 \qquad
 \kappa_N(f)\le B_L(d,C):=d(d+1)C^2+LC^4.
\]
For the specified finite-variance family, the raw target mass satisfies
\[
 \mu_{N,\gamma}(U_{N,f})\ge
 \exp\{-A_L(d,C)N^\gamma/2\}\qquad(\gamma\ge1).
\]
These are the compression, cloning, and target-cylinder bounds
of~\cite{strong,above,meanfield}, recorded with explicit scalar formulas
in~\cite{online}.

For strong prior, fix any $0<\beta<1$ and any admissible finite $N$.
There is $\sigma_0>0$ such that for every $0<\sigma<\sigma_0$, the class
\[
 \mathcal A=\{g\in\mathcal H_N:\kappa_N(g)\le\kappa_N(f)\}
\]
contains $f$ and satisfies
\[
 \log|\mathcal A|\le h:=F_L(d,B_L(d,C)),\qquad
 \pi(\mathcal A^c)\le\beta\pi(f).
\]
The threshold may depend on $N,d,L,f,\beta$ and the activation; no
width-uniform variance threshold is asserted. One threshold suffices for
all subsequent noiseless histories and the entire input stream.

For every fixed $\gamma>1$, put $A=A_L(d,C)$ and
\[
 B_\beta=4A+8\log(L/\beta),\qquad
 a_\beta=[8L(\sqrt{B_\beta}+1)^{L-2}]^{-1},
\]
\[
 T_\beta=\left\lceil\max\left\{
 C,[8(d+1)]^{1/(\gamma-1)},(8/a_\beta)^{2/(\gamma-1)}
 \right\}\right\rceil,
 \quad h_\beta=F_L(d,(L+3)B_\beta),
 \quad J_\beta=AT_\beta^\gamma/2.
\]
At $N\ge T_\beta$, the class
$\mathcal A=\{g:\kappa_N(g)\le(L+3)B_\beta\}$ contains $f$ and
satisfies $\log|\mathcal A|\le h_\beta$ and
$\pi(\mathcal A^c)\le\beta\pi(f)$.
At $C\le N<T_\beta$, instead $-\log\pi(f)\le J_\beta$.
Both budgets are polynomially bounded in $d,C,\log(1/\beta)$ at fixed
$L,\gamma$; their degrees may depend on $\gamma$.

For mean-field tanh, the functions $T_L,G_L$ have the following
structural meaning: for $B\ge4$ and $N\ge T_L(d,B)$, the raw total
mass of full-margin functions outside a finite class of log cardinality
at most $G_L(d,B)$ is at most $L e^{-NB/8}$.
For $\gamma=1$, use the same $A,B_\beta$ and set
\[
 T_\beta=\max\{C,T_L(d,B_\beta)\},\qquad
 h_\beta=G_L(d,B_\beta),\qquad J_\beta=AT_\beta/2.
\]
At $N\ge T_\beta$ there is a finite class $\mathcal A$ containing $f$
with $\log|\mathcal A|\le h_\beta$ and
$\pi(\mathcal A^c)\le\beta\pi(f)$; at $C\le N<T_\beta$,
$-\log\pi(f)\le J_\beta$. These budgets are polynomially bounded in
$d,C,\log(1/\beta)$ at fixed $L$.
\end{lemma}

\begin{proof}
The strong-prior assertion follows from
$2\sigma^2\log\mu_{N,\sigma}(U_{N,g})\to-\kappa_N(g)$ and
finiteness of $\mathcal H_N$: the total mass of functions with energy
strictly above the target's is $o(\pi(f))$. All energy ties stay in
$\mathcal A$; no assertion about equal weights is needed~\cite{strong,online}.
For $\gamma>1$, the spectral truncation estimate gives raw exceptional
mass at most $L\exp(-B_\beta N^\gamma/8)$ under the displayed width
conditions~\cite{above}. Dividing by the raw target mass gives at most
\[
 L\exp\{-(B_\beta/8-A/2)N^\gamma\}
 =L(\beta/L)^{N^\gamma}\le\beta.
\]
The mean-field finite-description theorem gives the identical calculation
with exceptional mass $L\exp(-NB_\beta/8)$ and width threshold
$T_L(d,B_\beta)$~\cite{meanfield}. In either case the resulting class
must contain $f$, since otherwise its relative outside mass is at least
one. At smaller widths,
$-\log\pi(f)\le-\log\mu(U_{N,f})\le AN^\gamma/2$.
The polynomial bounds and explicit quantities are also collected
in~\cite{online}. The classes and tail estimates are fixed before the
input process; exact conditioning therefore bounds posterior mass outside
each such class by $\beta$ after every realizable history.
\end{proof}

\begin{corollary}[Gaussian function ranks]
\label{cor:neural-rank}
For any prior covered by Lemma~\ref{lem:neural-inputs}, let
\[
 R_f=\#\{g\in\mathcal H_N:\pi(g)\ge\pi(f)\}.
\]
Then $\log R_f\le H_\beta$, where $H_\beta=h$ in the strong-prior
case and $H_\beta=\max\{h_\beta,J_\beta\}$ in either finite-variance
case. In particular, choosing $\beta=1/2$ gives a bound polynomial in
$d,C$ at the fixed parameters of the regime, uniformly for every
$N\ge C$. This rank is unchanged if raw global-margin masses replace
the conditioned function probabilities. The scalar $H_\beta$ bounds
both proof branches; it does not assert a common relevant class in the
smaller-width branch.
\end{corollary}

\begin{proof}
In a relevant-class branch, every $g\notin\mathcal A$ has
$\pi(g)\le\pi(\mathcal A^c)\le\beta\pi(f)<\pi(f)$.
Thus all functions counted by $R_f$, including every exact tie, belong
to $\mathcal A$, and $\log R_f\le\log|\mathcal A|$.
In a direct-mass branch, $R_f\pi(f)\le1$, so
$\log R_f\le-\log\pi(f)\le J_\beta$.
Global-margin conditioning divides all raw function masses by one common
normalizer and therefore preserves the ranking and its ties.
\end{proof}

When using strong-prior predictive limits, one takes limits of the
posterior laws after conditioning on each realizable history. This does
not require conditioning a limiting prior, which may assign the target
zero mass.
